% Options for packages loaded elsewhere
\PassOptionsToPackage{unicode}{hyperref}
\PassOptionsToPackage{hyphens}{url}
\documentclass[
  11pt,
  a4paper,
]{article}
\usepackage{xcolor}
\usepackage[top=2.0cm,bottom=2.0cm,left=2.0cm,right=2.0cm]{geometry}
\usepackage{amsmath,amssymb}
\setcounter{secnumdepth}{-\maxdimen} % remove section numbering
\usepackage{iftex}
\ifPDFTeX
  \usepackage[T1]{fontenc}
  \usepackage[utf8]{inputenc}
  \usepackage{textcomp} % provide euro and other symbols
\else % if luatex or xetex
  \usepackage{unicode-math} % this also loads fontspec
  \defaultfontfeatures{Scale=MatchLowercase}
  \defaultfontfeatures[\rmfamily]{Ligatures=TeX,Scale=1}
\fi
\usepackage{lmodern}
\ifPDFTeX\else
  % xetex/luatex font selection
  \setmainfont[]{Times New Roman}
  \setsansfont[]{Arial}
  \setmonofont[]{Menlo}
\fi
% Use upquote if available, for straight quotes in verbatim environments
\IfFileExists{upquote.sty}{\usepackage{upquote}}{}
\IfFileExists{microtype.sty}{% use microtype if available
  \usepackage[]{microtype}
  \UseMicrotypeSet[protrusion]{basicmath} % disable protrusion for tt fonts
}{}
\makeatletter
\@ifundefined{KOMAClassName}{% if non-KOMA class
  \IfFileExists{parskip.sty}{%
    \usepackage{parskip}
  }{% else
    \setlength{\parindent}{0pt}
    \setlength{\parskip}{6pt plus 2pt minus 1pt}}
}{% if KOMA class
  \KOMAoptions{parskip=half}}
\makeatother
\usepackage{longtable,booktabs,array}
\newcounter{none} % for unnumbered tables
\usepackage{calc} % for calculating minipage widths
% Correct order of tables after \paragraph or \subparagraph
\usepackage{etoolbox}
\makeatletter
\patchcmd\longtable{\par}{\if@noskipsec\mbox{}\fi\par}{}{}
\makeatother
% Allow footnotes in longtable head/foot
\IfFileExists{footnotehyper.sty}{\usepackage{footnotehyper}}{\usepackage{footnote}}
\makesavenoteenv{longtable}
\setlength{\emergencystretch}{3em} % prevent overfull lines
\providecommand{\tightlist}{%
  \setlength{\itemsep}{0pt}\setlength{\parskip}{0pt}}
% 1-column header for CystoDS section report
\PassOptionsToPackage{table,xcdraw}{xcolor}
\usepackage[section]{placeins}
\usepackage{caption}
\usepackage{fancyhdr}
\usepackage{titlesec}
\usepackage{enumitem}
\usepackage{seqsplit}
\usepackage{booktabs}
\usepackage{graphicx}
\usepackage{microtype}
\usepackage{tabularx}
\usepackage{adjustbox}
\usepackage{float}

\definecolor{CystoBlue}{HTML}{1F4E79}
\definecolor{CystoTeal}{HTML}{2A7F62}
\definecolor{CystoGray}{HTML}{4A5568}
\definecolor{CystoDark}{HTML}{1A202C}

\graphicspath{{../../Docs/}{Docs/}{paper_assets/}{../paper_assets/}}

\captionsetup{font=small,labelfont={bf,color=CystoBlue}}
\captionsetup[table]{position=top,skip=3pt}
\captionsetup[figure]{position=bottom,skip=4pt}

\renewcommand{\figurename}{Hình}
\renewcommand{\tablename}{Bảng}

\titleformat{\section}{\Large\bfseries\color{CystoBlue}}{\thesection.}{0.4em}{}
\titleformat{\subsection}{\large\bfseries\color{CystoTeal}}{\thesubsection.}{0.4em}{}
\titleformat{\subsubsection}{\normalsize\bfseries\color{CystoDark}}{\thesubsubsection.}{0.4em}{}

\setlength{\parindent}{1.0em}
\setlength{\parskip}{0.35em plus 0.05em minus 0.05em}
\setlist{nosep,leftmargin=1.5em,topsep=0.3em}

\pagestyle{fancy}
\fancyhf{}
\fancyhead[L]{\small\color{CystoGray}CystoDS -- Báo cáo Tập Dữ Liệu và Giao Thức (Stage 00)}
\fancyhead[R]{\small\color{CystoGray}CystoDS 2026}
\fancyfoot[C]{\small\color{CystoGray}\thepage}
\setlength{\headheight}{14pt}
\renewcommand{\headrulewidth}{0.35pt}

\AtBeginDocument{%
  \hypersetup{
    colorlinks=true,
    linkcolor=CystoBlue,
    urlcolor=CystoTeal,
    citecolor=CystoBlue,
  }%
}
\usepackage{bookmark}
\IfFileExists{xurl.sty}{\usepackage{xurl}}{} % add URL line breaks if available
\urlstyle{same}
\hypersetup{
  pdftitle={CystoDS: Báo cáo Tập Dữ Liệu{,} Bối Cảnh và Giao Thức Đánh Giá Độc Lập Bệnh Nhân},
  hidelinks,
  pdfcreator={LaTeX via pandoc}}

\title{CystoDS: Báo cáo Tập Dữ Liệu, Bối Cảnh và Giao Thức Đánh Giá Độc
Lập Bệnh Nhân}
\usepackage{etoolbox}
\makeatletter
\providecommand{\subtitle}[1]{% add subtitle to \maketitle
  \apptocmd{\@title}{\par {\large #1 \par}}{}{}
}
\makeatother
\subtitle{Kiểm định 8.067 ảnh, 160 bệnh nhân, phân tích liên hệ nghiên
cứu và 3 phân hoạch hold-out chuẩn hóa (Stage 00)}
\author{}
\date{Báo cáo nghiên cứu 2026}

\begin{document}
\maketitle

\section{CystoDS: Báo cáo Tập Dữ Liệu, Bối Cảnh và Giao Thức Đánh Giá
Độc Lập Bệnh
Nhân}\label{cystods-buxe1o-cuxe1o-tux1eadp-dux1eef-liux1ec7u-bux1ed1i-cux1ea3nh-vuxe0-giao-thux1ee9c-ux111uxe1nh-giuxe1-ux111ux1ed9c-lux1eadp-bux1ec7nh-nhuxe2n}

\subsection{Kiểm định 8.067 ảnh, 160 bệnh nhân, phân tích liên hệ nghiên
cứu và 3 phân hoạch hold-out chuẩn hóa (Stage
00)}\label{kiux1ec3m-ux111ux1ecbnh-8.067-ux1ea3nh-160-bux1ec7nh-nhuxe2n-phuxe2n-tuxedch-liuxean-hux1ec7-nghiuxean-cux1ee9u-vuxe0-3-phuxe2n-houx1ea1ch-hold-out-chuux1ea9n-huxf3a-stage-00}

\subsection{1. Đặt vấn đề \& Đóng góp khoa học (Introduction \&
Contributions)}\label{ux111ux1eb7t-vux1ea5n-ux111ux1ec1-ux111uxf3ng-guxf3p-khoa-hux1ecdc-introduction-contributions}

\subsubsection{1.1. Bối cảnh lâm sàng \& Tính chất dữ liệu nội soi bàng
quang}\label{bux1ed1i-cux1ea3nh-luxe2m-suxe0ng-tuxednh-chux1ea5t-dux1eef-liux1ec7u-nux1ed9i-soi-buxe0ng-quang}

\begin{itemize}
\tightlist
\item
  \textbf{Ý nghĩa lâm sàng:} Nội soi bàng quang (Cystoscopy) là tiêu
  chuẩn vàng chẩn đoán ung thư biểu mô đường niệu; tuy nhiên, thị trường
  lâm sàng đối mặt với tỷ lệ bỏ sót tổn thương phẳng (CIS) và cảnh báo
  giả do niêm mạc viêm hoặc mốc giải phẫu.
\item
  \textbf{Tính đa tầng tự nhiên của chẩn đoán (Diagnostic Hierarchy):}

  \begin{itemize}
  \tightlist
  \item
    Mức 1: Định vị nhanh vùng tổn thương nghi ngờ (ROI Detection: Tổn
    thương vs Bình thường).
  \item
    Mức 2: Phân nhóm bệnh cảnh thô (Coarse Grouping: Ác tính, Không ác
    tính, Niêm mạc lành, Mốc giải phẫu, Dị vật/Dụng cụ).
  \item
    Mức 3: Phân loại chẩn đoán chi tiết (Fine-Grained Diagnostic
    Categories: 22 phân lớp chẩn đoán bao gồm tổn thương ác tính, lành
    tính, mốc giải phẫu và dụng cụ can thiệp).
  \end{itemize}
\item
  \textbf{Thách thức mất cân bằng dài đuôi cực đoan ở cấp độ bệnh nhân:}

  \begin{itemize}
  \tightlist
  \item
    Đỉnh phân phối (Head): Niêm mạc bình thường chiếm áp đảo 79,16\%
    tổng số ảnh.
  \item
    Đuôi phân phối (Tail): Các tổn thương tiền ung thư
    (\texttt{PreMalignant}), ung thư biểu mô tại chỗ (\texttt{CIS}), nhú
    niệu mạc (\texttt{UrothelialPapilloma}) chỉ xuất hiện ở vài bệnh
    nhân cá biệt.
  \end{itemize}
\end{itemize}

\subsubsection{1.2. Mâu thuẫn kỹ thuật cốt lõi trong học đa nhiệm phân
cấp}\label{muxe2u-thuux1eabn-kux1ef9-thuux1eadt-cux1ed1t-luxf5i-trong-hux1ecdc-ux111a-nhiux1ec7m-phuxe2n-cux1ea5p}

\begin{itemize}
\tightlist
\item
  \textbf{Mâu thuẫn 1 -- Xung đột Gradient biểu diễn vs Cân bằng đuôi
  dài (Representation vs Classifier Trade-off):}

  \begin{itemize}
  \tightlist
  \item
    Huấn luyện 1 giai đoạn chung (1-Stage Joint) với các hàm loss đuôi
    dài mạnh làm méo mó không gian đặc trưng tổng quát, dẫn đến suy giảm
    độ chính xác lớp đầu (Coarse/Binary).
  \end{itemize}
\item
  \textbf{Mâu thuẫn 2 -- Thắt cổ chai phân cấp sớm (Early Hierarchy
  Bottleneck):}

  \begin{itemize}
  \tightlist
  \item
    Ràng buộc phân cấp phả hệ (\(L_{\text{hierarchy}}\)) nếu áp đặt cứng
    ngay từ epoch 1 sẽ gò bó gradient của Backbone, cản trở việc học các
    đặc trưng thị giác cơ bản.
  \end{itemize}
\item
  \textbf{Mâu thuẫn 3 -- Nguy cơ trôi dạt ranh giới khi tinh chỉnh phân
  tách:}

  \begin{itemize}
  \tightlist
  \item
    Nếu chỉ nắn Fine Head ở giai đoạn sau mà không có cơ chế cô lập tham
    số, ranh giới quyết định của Binary và Coarse sẽ bị trôi dạt.
  \end{itemize}
\end{itemize}

\subsubsection{1.3. Ba đóng góp cốt lõi của nghiên
cứu}\label{ba-ux111uxf3ng-guxf3p-cux1ed1t-luxf5i-cux1ee7a-nghiuxean-cux1ee9u}

\begin{enumerate}
\def\labelenumi{\arabic{enumi}.}
\tightlist
\item
  \textbf{Khung học phân cấp nhận biết phân bố bệnh nhân (Patient-Aware
  Hierarchical Learning):} Tích hợp kiến trúc Shared Late-Stage
  Swin-Tiny với hàm mất mát \textbf{Patient-Prior Smoothed Balanced
  Softmax (PP-SBS)} điều hòa theo căn bậc hai số lượng bệnh nhân thực tế
  \(\pi_j = (\text{patients}_j + \epsilon)^{0{,}5}\) và Supervised
  Contrastive Loss (\(L_{\text{supcon}}\)).
\item
  \textbf{Chiến lược tối ưu hóa tuần tự tách rời ba giai đoạn
  (Three-Stage Decoupled Optimization):} Tách bạch rõ ràng:
  \emph{Representation Learning (với Curriculum Warmup
  \(w_{\text{hrc}}(t)\))} \rightarrow *Coarse Alignment (với cơ chế cô
  lập tham số)* \rightarrow *Fine Alignment*.
\item
  \textbf{Cơ chế suy luận phả hệ tăng cường và kiểm định thực nghiệm
  minh bạch:} Tích lũy xác suất từ 22 phân lớp chi tiết về 5 nhóm cha
  (\(P_{\text{from\_fine}}\)) kết hợp hòa trộn (\(\lambda=0{,}25\)),
  tăng Coarse Accuracy thêm \(+5{,}24\%\) trên tập Test độc lập. Báo cáo
  đầy đủ khoảng tin cậy 95\%, kiểm định thống kê cặp và phân tích độ
  nhạy Grad-CAM.
\end{enumerate}

\begin{center}\rule{0.5\linewidth}{0.5pt}\end{center}

\subsection{2. Công trình liên quan \& Định vị nghiên cứu (Related
Work)}\label{cuxf4ng-truxecnh-liuxean-quan-ux111ux1ecbnh-vux1ecb-nghiuxean-cux1ee9u-related-work}

\subsubsection{2.1. Phân loại ảnh nội soi bàng quang bằng Deep
Learning}\label{phuxe2n-loux1ea1i-ux1ea3nh-nux1ed9i-soi-buxe0ng-quang-bux1eb1ng-deep-learning}

\begin{itemize}
\tightlist
\item
  \textbf{Nghiên cứu gốc CystoDS (Lee et al., 2026 {[}1{]}):} Khảo sát
  ResNet, ResNeXt, HRNet, Swin-Transformer cho bài toán nhị phân ROI;
  chỉ dừng lại ở phân loại phẳng trên 1 split duy nhất.
\item
  \textbf{Các nghiên cứu quốc tế liên quan:}

  \begin{itemize}
  \tightlist
  \item
    \emph{Shkolyar et al.~(2019 {[}5{]}):} Phát hiện khối u trên video
    WLC (Sensitivity 90,9\%, Specificity 98,6\%), không phân loại chi
    tiết đa tầng.
  \item
    \emph{Wu et al.~(2022 {[}6{]}):} Đa trung tâm 69.204 ảnh, phân loại
    nhị phân ung thư / lành tính.
  \item
    \emph{Lazo et al.~(2023 {[}7{]}):} Phân loại 4 nhóm mô bán giám sát
    trên WLI/NBI (1.754 ảnh/23 BN).
  \item
    \emph{Abd El-Aziz et al.~(2025 {[}9{]}):} EfficientNet-B3 phân loại
    4 lớp trên bộ dữ liệu EBTC.
  \item
    \emph{Wang et al.~(2026 {[}10{]}):} Chẩn đoán NBI đa trung tâm, kết
    hợp phân độ mô học.
  \end{itemize}
\end{itemize}

\subsubsection{2.2. Học phân cấp \& Cân bằng phân phối đuôi
dài}\label{hux1ecdc-phuxe2n-cux1ea5p-cuxe2n-bux1eb1ng-phuxe2n-phux1ed1i-ux111uuxf4i-duxe0i}

\begin{itemize}
\tightlist
\item
  \textbf{Hierarchical Classification:} Tận dụng cây phả hệ y học để
  phạt các lỗi sai nhánh nghiêm trọng hơn lỗi sai trong cùng nhóm cha.
\item
  \textbf{Decoupled Representation \& Classifier Learning (Kang et al.,
  2020; cRT):} Tách rời việc học đặc trưng trên phân phối tự nhiên và
  cân bằng lại phân loại; được mở rộng trong nghiên cứu này thành cơ chế
  3 giai đoạn tuần tự.
\item
  \textbf{Balanced Meta-Softmax (Ren et al., 2020 {[}3{]}):} Dịch chuyển
  biên quyết định theo phân phối mẫu; được chúng tôi cải tiến thành
  \textbf{Patient-Prior Smoothed Balanced Softmax (PP-SBS)} với prior
  mượt theo số lượng bệnh nhân
  (\(\text{prior}_j = (\text{patients}_j + \epsilon)^{0{,}5}\)).
\end{itemize}

\subsubsection{2.3. Bảng đối chiếu tổng hợp các công trình liên
quan}\label{bux1ea3ng-ux111ux1ed1i-chiux1ebfu-tux1ed5ng-hux1ee3p-cuxe1c-cuxf4ng-truxecnh-liuxean-quan}

\begin{table}[H]
\centering
\caption{Bảng 1.}
\label{tab:sec_table1}
\vspace{3pt}
\adjustbox{max width=\textwidth}{
\begin{tabular}{lrrrrr}
\toprule
Nghiên cứu & Bộ dữ liệu / Cỡ mẫu & Số lượng Lớp / Tầng & Độc lập Bệnh nhân & Phương pháp Cốt lõi & Điểm hạn chế chính \\
\midrule
\textbf{Lee et al.~(CystoDS gốc) {[}1{]}} & 8.067 ảnh / 160 BN & Nhị
phân ROI (2 lớp) & 1 Split riêng & Swin-Transformer Large & Chỉ làm nhị
phân, chưa khai thác 5 coarse / 22 fine \\
\textbf{Shkolyar et al.~{[}5{]}} & Video WLC nội bộ & ROI Detection & Có
& CNN Object Detection & Không phân loại chẩn đoán đa tầng \\
\textbf{Wu et al.~{[}6{]}} & 69.204 ảnh / 10.729 BN & Ung thư vs Lành
tính & Đa trung tâm & CNN Classifier & Nhị phân phẳng, thiếu chi tiết
tổn thương \\
\textbf{Lazo et al.~{[}7{]}} & 1.754 ảnh / 23 BN & 4 lớp phẳng & Có &
Semi-supervised & Cỡ mẫu nhỏ, không có cấu trúc phả hệ \\
\textbf{Wang et al.~{[}10{]}} & NBI đa trung tâm & Phân độ ung thư & Có
& Multitask NBI & Phụ thuộc ánh sáng dải hẹp NBI \\
\textbf{Nghiên cứu này (CystoHier)} & \textbf{8.067 ảnh / 160 BN} &
\textbf{3 tầng (2 / 5 / 22 lớp)} & \textbf{3 Splits chuẩn hóa} &
\textbf{CystoHier (3S + Ens.)} & \textbf{Giải quyết đồng thời 3 tầng
nhãn và đuôi dài} \\
\bottomrule
\end{tabular}
}
\end{table}

\begin{center}\rule{0.5\linewidth}{0.5pt}\end{center}

\subsection{3. Kiểm toán Dữ liệu \& Giao thức Nghiên cứu (Stage 00
Protocol)}\label{kiux1ec3m-touxe1n-dux1eef-liux1ec7u-giao-thux1ee9c-nghiuxean-cux1ee9u-stage-00-protocol}

\subsubsection{3.1. Thống kê bộ dữ liệu CystoDS qua 3 tầng phân
cấp}\label{thux1ed1ng-kuxea-bux1ed9-dux1eef-liux1ec7u-cystods-qua-3-tux1ea7ng-phuxe2n-cux1ea5p}

\begin{table}[H]
\centering
\caption{Bảng 2.}
\label{tab:sec_table2}
\vspace{3pt}
\adjustbox{max width=\textwidth}{
\begin{tabular}{lrrr}
\toprule
Tầng nhãn & Số lớp & Chi tiết phân nhóm lâm sàng & Số lượng mẫu (Raw) \\
\midrule
\textbf{Layer 1: Binary} & 2 & \textbf{ROI} (Tổn thương) /
\textbf{Non-ROI} (Bình thường \& Dị vật) & 1.219 ROI / 6.848 Non-ROI \\
\textbf{Layer 2: Coarse} & 5 & Malignant (998), Non-malignant (221),
Normal mucosa (6.386), Landmarks (211), Foreign bodies (251) & 8.067 \\
\textbf{Layer 3: Fine} & 22 & 4 Ác tính, 8 Lành tính/Viêm, 6 Mốc giải
phẫu, 4 Dị vật/Dụng cụ & 8.067 (Long-tailed) \\
\bottomrule
\end{tabular}
}
\end{table}

\subsubsection{3.2. Kiểm định rò rỉ dữ liệu \& Kiểm soát thiên lệch niêm
mạc}\label{kiux1ec3m-ux111ux1ecbnh-ruxf2-rux1ec9-dux1eef-liux1ec7u-kiux1ec3m-souxe1t-thiuxean-lux1ec7ch-niuxeam-mux1ea1c}

\begin{itemize}
\tightlist
\item
  \textbf{Kiểm tra trùng lặp (Image Duplication):} 0 cặp ảnh trùng hash
  SHA-256 trong toàn bộ 8.067 ảnh.
\item
  \textbf{Giới hạn niêm mạc bình thường
  (\texttt{normal\_mucosa\_limit:\ 540}):} Giới hạn tối đa 540 ảnh
  Normal mucosa trong tập Train để tránh làm loãng đặc trưng bệnh lý.
\item
  \textbf{Giao thức 3-Fold Patient-Disjoint Holdout:}

  \begin{itemize}
  \tightlist
  \item
    Tỉ lệ phân chia: 70\% Train (\textasciitilde1.550 ảnh / 112 BN),
    15\% Validation (\textasciitilde330 ảnh / 24 BN), 15\% Test
    (\textasciitilde335 ảnh / 24 BN).
  \item
    Khóa toàn bộ danh tính bệnh nhân giữa các tập để đảm bảo độ trung
    thực lâm sàng tuyệt đối.
  \end{itemize}
\end{itemize}

\begin{center}\rule{0.5\linewidth}{0.5pt}\end{center}

\end{document}
